% Part: incompleteness
% Chapter: incompleteness-provability
% Section: rosser-thm

\documentclass[../../../include/open-logic-section]{subfiles}

\begin{document}

\olfileid{inc}{inp}{ros}

\olsection{قضیهٔ راسر}

آیا می‌توانیم اثبات گودل را چنان اصلاح کنیم که به نتیجه‌ای قوی‌تر برسیم
و «$\omega$-سازگار» را صرفاً با «سازگار» جایگزین کنیم؟ پاسخ با استفاده
از ترفندی که راسر کشف کرد مثبت است. ترفند راسر استفاده از محمول
!!{derivability} «اصلاح‌شده»ای چون~$\ORProv_T(y)$ به جای
$\OProv[\Th{T}](y)$ است.

\begin{thm}
\ollabel{thm:rosser}
فرض کنید~$\Th{T}$ نظریه‌ای دلخواه، سازگار و !!{axiomatizable} و
گسترشی از~$\Th{Q}$ باشد. آنگاه~$\Th{T}$ کامل نیست.
\end{thm}

\begin{proof}
به یاد آورید که~$\OProv[\Th{T}](y)$ به صورت
$\lexists[x][\OPrf[\Th{T}](x, y)]$ تعریف می‌شود، که در آن
$\OPrf[\Th{T}](x, y)$ رابطهٔ تصمیم‌پذیری را بازنمایی می‌کند که برقرار
است اگر و تنها اگر~$x$ عدد گودل !!a{derivation}ی از !!{sentence}ی با
عدد گودل~$y$ باشد. رابطه‌ای که میان~$x$ و~$y$ برقرار است هرگاه~$x$
عدد گودل یک \emph{ابطال} جملهٔ دارای عدد گودل~$y$ باشد نیز تصمیم‌پذیر
است. بگذارید~$\fn{not}(x)$ تابع بازگشتی اولیه‌ای باشد که چنین عمل
می‌کند: اگر~$x$ رمز فرمول~$!A$ باشد، $\fn{not}(x)$ رمزی برای
$\lnot !A$ است. آنگاه~$\Refut[\Th{T}](x, y)$ برقرار است اگر و تنها اگر
$\Prf[\Th{T}](x, \fn{not}(y))$. بگذارید~$\ORefut[\Th{T}](x, y)$ آن را
بازنمایی کند. پس اگر~$\Th{T} \Proves \lnot !A$ و~$\delta$
!!{derivation} متناظر باشد، آنگاه
$\Th{Q} \Proves \ORefut[\Th{T}](\gn{\delta}, \gn{!A})$. اکنون
$\ORProv[\Th{T}](y)$ را چنین تعریف می‌کنیم:
\[
\lexists[x][(\OPrf[\Th{T}](x,y) \land \lforall[z][(z < x \lif \lnot
  \ORefut[\Th{T}](z,y))])].
\]
به‌طور تقریبی، $\ORProv[\Th{T}](y)$ می‌گوید «اثباتی برای~$y$ در
$\Th{T}$ وجود دارد و هیچ ابطال کوتاه‌تری برای~$y$ وجود ندارد.» با فرض
سازگاری~$\Th{T}$، $\ORProv[\Th{T}](y)$ دربارهٔ همان اعدادی صادق است که
$\OProv[\Th{T}](y)$ دربارهٔ آن‌ها صادق است؛ اما از دیدگاه
\emph{اثبات‌پذیری} در~$\Th{T}$ (و اکنون می‌دانیم که صدق و اثبات‌پذیری
متفاوت‌اند)، این دو ویژگی‌های متفاوتی دارند. اگر~$\Th{T}$
\emph{نا}سازگار باشد، این دو دربارهٔ اعداد یکسانی برقرار \emph{نیستند}!{}
($\ORProv[\Th{T}](y)$ را اغلب چنین می‌خوانند: «$y$ به شیوهٔ راسر
اثبات‌پذیر است.» اما چون، چنان‌که دیدیم، اثبات‌پذیری راسر نوع ویژه‌ای از
اثبات‌پذیری نیست---در نظریه‌های ناسازگار !!{sentence}sی وجود دارند که
اثبات‌پذیرند اما به شیوهٔ راسر اثبات‌پذیر نیستند---این خوانش ممکن است
گمراه‌کننده باشد. برای پرهیز از این ابهام می‌توانید آن را چنین بخوانید:
«$y$ شِماثبات‌پذیر است».)

بنا بر لم نقطهٔ ثابت، فرمولی چون~$!R_\Th{T}$ وجود دارد که
\begin{equation}
  \Th{Q} \Proves !R_\Th{T} \liff \lnot \ORProv[\Th{T}](\gn{!R_\Th{T}}).
  \ollabel{RT}
\end{equation}
برخلاف اثبات~\olref[1in]{thm:first-incompleteness}، اینجا ادعا می‌کنیم
اگر~$\Th{T}$ سازگار باشد، $\Th{T}$، $!R_\Th{T}$ را !!{derive}
نمی‌کند و~$\Th{T}$، $\lnot !R_\Th{T}$ را نیز !!{derive} نمی‌کند. (به
بیان دیگر، به فرض $\omega$-سازگاری نیاز نداریم.)

نخست نشان می‌دهیم~$\Th{T} \Proves/ !R_{\Th{T}}$. خلاف آن را فرض کنید؛
پس !!a{derivation}ی از~$!R_{\Th{T}}$ بر پایهٔ~$\Th{T}$ وجود دارد.
بگذارید~$n$ عدد گودل آن باشد. آنگاه
$\Th{Q} \Proves \OPrf[\Th{T}](\num{n}, \gn{!R_{\Th{T}}})$، زیرا
$\OPrf[\Th{T}]$ رابطهٔ~$\Prf[\Th{T}]$ را در~$\Th{Q}$ بازنمایی می‌کند.
همچنین برای هر~$k < n$، عدد~$k$ عدد گودل !!a{derivation}ی از
$\lnot !R_{\Th{T}}$ نیست، زیرا~$\Th{T}$ سازگار است. پس برای هر
$k < n$ داریم
$\Th{Q} \Proves \lnot \ORefut[\Th{T}](\num{k}, \gn{!R_{\Th{T}}})$.
بنا بر~\olref[req][min]{lem:less-nsucc}،
$\Th{Q} \Proves \lforall[z][(z < \num{n} \lif \lnot \ORefut[\Th{T}](z,
  \gn{!R_{\Th{T}}}))]$. بنابراین
\[
\Th{Q} \Proves \lexists[x][(\OPrf[\Th{T}](x,\gn{!R_{\Th{T}}}) \land \lforall[z][(z
    < x \lif \lnot \ORefut[\Th{T}](z,\gn{!R_{\Th{T}}}))])],
\]
اما این همان~$\ORProv[\Th{T}](\gn{!R_{\Th{T}}})$ است. بنا بر
\olref{RT}، $\Th{Q} \Proves \lnot !R_{\Th{T}}$. چون~$\Th{T}$ گسترشی
از~$\Th{Q}$ است، $\Th{T} \Proves \lnot !R_{\Th{T}}$ نیز. فرض کرده
بودیم~$\Th{T} \Proves !R_{\Th{T}}$، پس~$\Th{T}$ برخلاف فرض قضیه
ناسازگار می‌شد.

اکنون نشان می‌دهیم~$\Th{T} \Proves/ \lnot !R_{\Th{T}}$. باز هم خلاف
آن را فرض کنید و بگذارید~$n$ عدد گودل !!a{derivation}ی از
$\lnot !R_{\Th{T}}$ باشد. آنگاه~$\Refut[\Th{T}](n, \Gn{!R_{\Th{T}}})$
برقرار است، و چون~$\ORefut[\Th{T}]$ رابطهٔ~$\Refut[\Th{T}]$ را در
$\Th{Q}$ بازنمایی می‌کند،
$\Th{Q} \Proves \ORefut[\Th{T}](\num{n}, \gn{!R_{\Th{T}}})$.
دوباره نشان می‌دهیم در این صورت~$\Th{T}$ ناسازگار خواهد بود، زیرا
$!R_{\Th{T}}$ را نیز !!{derive} می‌کند. چون
\begin{align*}
\Th{Q} & \Proves !R_{\Th{T}} \liff \lnot \ORProv[\Th{T}](\gn{!R_{\Th{T}}}),
\intertext{و چون~$\Th{T}$ گسترشی از~$\Th{Q}$ است، کافی است نشان دهیم}
\Th{Q} & \Proves \lnot \ORProv[\Th{T}](\gn{!R_{\Th{T}}}).
\end{align*}
!!{sentence}~$\lnot
\ORProv[\Th{T}](\gn{!R_{\Th{T}}})$، یعنی
\begin{align*}
  \lnot & \lexists[x][(\OPrf[\Th{T}](x,\gn{!R_{\Th{T}}}) \land
    \lforall[z][(z < x \lif \lnot \ORefut[\Th{T}](z,\gn{!R_{\Th{T}}}))])],
  \intertext{از نظر منطقی هم‌ارز است با}
  & \lforall[x][(\OPrf[\Th{T}](x,\gn{!R_{\Th{T}}}) \lif
    \lexists[z][(z < x \land \ORefut[\Th{T}](z,\gn{!R_{\Th{T}}}))])].
\end{align*}
با استفاده از منطق و واقعیت‌هایی دربارهٔ آنچه~$\Th{Q}$ !!{derive}s می‌کند،
به‌طور غیرصوری استدلال می‌کنیم. فرض کنید~$x$ دلخواه باشد و
$\OPrf[\Th{T}](x, \gn{!R_{\Th{T}}})$. از پیش می‌دانیم
$\Th{T} \Proves/ !R_{\Th{T}}$؛ پس برای هر~$k$ داریم
$\Th{Q} \Proves \lnot \OPrf[\Th{T}](\num{k}, \gn{!R_{\Th{T}}})$.
بنابراین برای هر~$k$ نتیجه می‌شود~$\eq/[x][\num{k}]$. به‌ویژه،
(الف)~$\eq/[x][\num{n}]$. همچنین
$\lnot(\eq[x][\num{0}] \lor \eq[x][\num{1}] \lor \dots \lor
\eq[x][\num{n-1}])$ را داریم؛ پس بنا بر
\olref[req][min]{lem:less-nsucc}، (ب)~$\lnot(x < \num{n})$. بنا بر
\olref[req][min]{lem:trichotomy}، $\num{n} < x$. چون
$\Th{Q} \Proves \ORefut[\Th{T}](\num{n}, \gn{!R_{\Th{T}}})$، داریم
$\num{n} < x \land \ORefut[\Th{T}](\num{n}, \gn{!R_{\Th{T}}})$، و از آن
$\lexists[z][(z < x \land \ORefut[\Th{T}](z,\gn{!R_{\Th{T}}}))]$.
چون~$x$ دلخواه بود، چنان‌که لازم بود به دست می‌آوریم:
\[
\lforall[x][(\OPrf[\Th{T}](x,\gn{!R_{\Th{T}}}) \lif
  \lexists[z][(z < x \land \ORefut[\Th{T}](z,\gn{!R_{\Th{T}}}))])].
\]
\end{proof}

\begin{prob}
دو مجموعهٔ~$A$ و~$B$ از اعداد طبیعی را به‌طور محاسبه‌پذیر
\emph{جدایی‌ناپذیر} می‌نامیم اگر هیچ مجموعهٔ !!{decidable}~$X$ای وجود
نداشته باشد که~$A \subseteq X$ و~$B \subseteq \Complement{X}$
($\Complement{X}$ نماد متمم است: $\Nat \setminus X$، که متمم~$X$
است). فرض کنید
$\Th{T}$ گسترشی سازگار و !!{axiomatizable} از~$\Th{Q}$ باشد. فرض کنید
$A$ مجموعهٔ اعداد گودل !!{sentence}sی اثبات‌پذیر در~$\Th{T}$ و~$B$
مجموعهٔ اعداد گودل جمله‌های ابطال‌پذیر در~$\Th{T}$ باشد. ثابت کنید
$A$ و~$B$ به‌طور محاسبه‌پذیر جدایی‌ناپذیرند.
\end{prob}

\end{document}
